% Bagian: sets-functions-relations
% Bab: size-of-sets
% Subbagian: enumerations-alt

\documentclass[../../../include/open-logic-section]{subfiles}

\begin{document}

\olfileid[id]{sfr}{siz}{enm-alt}

\olsection{Enumerasi dan Himpunan \usetoken{S}{enumerable}}

\begin{editorial}
  Bagian ini mendefinisikan enumerasi sebagai bijeksi dengan (segmen awal)
  $\Nat$, sebagaimana lazim dilakukan dalam teori himpunan. Karena itu,
  definisinya sedikit bertentangan dengan definisi dalam \olref[enm]{sec}, dan
  bagian ini mengulang semua contoh di sana. Bagian ini juga sedikit lebih
  ringkas daripada bagian tersebut.
\end{editorial}

Kita dapat menentukan suatu himpunan hingga dengan mengenumerasi
!!{element}snya. Inilah yang kita lakukan ketika mendefinisikan himpunan
sebagai berikut:
\[
  A = \{a_1, a_2, \ldots, a_n\}.
\]
Dengan mengandaikan bahwa !!{element}s $a_1$, \dots, $a_n$ semuanya berbeda,
hal ini memberi kita !!a{bijection} antara $A$ dan $n$ bilangan asli pertama,
yakni $0$, \dots, $n-1$. Sebaliknya, karena setiap himpunan hingga hanya
mempunyai sejumlah hingga !!{element}s, setiap himpunan hingga dapat
ditempatkan dalam korespondensi semacam itu. Dengan kata lain, jika $A$
hingga, terdapat !!a{bijection} antara $A$ dan $\{0, \dots, n-1\}$, dengan $n$
sebagai banyaknya !!{element}s dari~$A$.

Jika kita mengizinkan jenis tertentu dari himpunan tak hingga, kita juga akan
mengizinkan beberapa himpunan tak hingga untuk dienumerasi. Kita dapat
menyatakannya secara saksama dengan mengatakan bahwa suatu himpunan tak hingga
dienumerasi oleh !!a{bijection} antara himpunan itu dan seluruh~$\Nat$.

\begin{defn}[Enumerasi, secara teori-himpunan]
Suatu \emph{enumerasi} himpunan $A$ adalah !!a{bijection} yang daerah hasilnya
$A$ dan domainnya berupa suatu himpunan awal bilangan asli $\{0, 1, \ldots,
n\}$ {atau} seluruh himpunan bilangan asli~$\Nat$.
\end{defn}

\begin{explain}
Penggunaan kata \emph{enumerasi} ini mempunyai landasan intuitif. Mengatakan
bahwa kita telah mengenumerasi suatu himpunan $A$ berarti mengatakan bahwa
terdapat !!a{bijection} $f$ yang memungkinkan kita menghitung anggota-anggota
himpunan $A$. Anggota ke-$0$ adalah $f(0)$, yang ke-1 adalah $f(1)$, \ldots,
yang ke-$n$ adalah $f(n)$\ldots.\footnote{Ya, kita menghitung mulai dari $0$.
Tentu saja kita juga dapat memulai dari~$1$. Ini tidak akan menimbulkan
perbedaan besar. Kita hanya perlu mengganti~$\Nat$ dengan~$\PosInt$.} Alasan
bagi hal ini dapat dibuat lebih jelas lagi dengan menambahkan definisi
berikut:
\end{explain}

\begin{defn}
  \ollabel{defn:enumerable}
  Suatu himpunan~$A$ adalah !!{enumerable} jika dan hanya jika $A =
  \emptyset$ atau terdapat suatu enumerasi dari~$A$. Kita mengatakan bahwa
  $A$ !!{nonenumerable} jika dan hanya jika $A$ tidak !!{enumerable}.
\end{defn}

\begin{explain}
Jadi, suatu himpunan !!{enumerable} jika dan hanya jika himpunan itu kosong
atau Anda dapat menggunakan suatu enumerasi untuk menghitung semua
!!{element}snya.
\end{explain}

\begin{ex}
Fungsi yang mengenumerasi bilangan asli cukup berupa fungsi identitas
$\Id{\Nat} \colon \Nat \to \Nat$ yang diberikan oleh $\Id{\Nat}(n) = n$.
Fungsi yang mengenumerasi bilangan asli \emph{positif}, $\Nat^+ = \Nat
\setminus \{0\}$, adalah fungsi $g(n) = n + 1$, yakni fungsi penerus.
\end{ex}

\begin{prob}
Tunjukkan bahwa suatu himpunan $A$ adalah !!{enumerable} jika dan hanya jika
$A = \emptyset$ atau terdapat !!a{surjection} $f\colon \Nat \to A$. Tunjukkan
bahwa $A$ adalah !!{enumerable} jika dan hanya jika terdapat !!a{injection}
$g\colon A \to \Nat$.
\end{prob}

\begin{ex}
Fungsi $f\colon \Nat \to \Nat$ dan $g \colon \Nat \to \Nat$ yang diberikan
oleh
\begin{align*}
  f(n) & = 2n \text{ dan}\\
  g(n) & = 2n+1
\end{align*}
masing-masing mengenumerasi bilangan asli genap dan bilangan asli ganjil.
Namun, tidak satu pun !!{surjective}, sehingga tidak satu pun merupakan
enumerasi dari $\Nat$.
\end{ex}

\begin{prob}
Definisikan suatu enumerasi dari bilangan kuadrat $1$, $4$, $9$, $16$, \dots
\end{prob}

\begin{ex}
Misalkan $\lceil x \rceil$ adalah fungsi \emph{pembulatan ke atas}, yang
membulatkan $x$ ke bilangan bulat terdekat di atasnya. Maka fungsi $f \colon
\Nat \to \Int$ yang diberikan oleh:
\[
  f(n) = (-1)^{n} \left\lceil\tfrac{n}{2}\right\rceil
\]
mengenumerasi himpunan bilangan bulat~$\Int$ sebagai berikut:
\[
\begin{array}{c c c c c c c c}
f(0) & f(1) & f(2) & f(3) & f(4) & f(5) & f(6) & \dots \\ \\
\big\lceil \tfrac{0}{2} \big\rceil & -\big\lceil \tfrac{1}{2}\big\rceil &  \big\lceil \tfrac{2}{2} \big\rceil & -\big\lceil \tfrac{3}{2} \big\rceil & \big\lceil \tfrac{4}{2} \big\rceil  & -\big\lceil \tfrac{5}{2}\big\rceil & \big\lceil \tfrac{6}{2} \big\rceil & \dots \\ \\
0 & -1 & 1 & -2 & 2 & -3 & 3& \dots
\end{array}
\]
Perhatikan bagaimana $f$ menghasilkan nilai-nilai $\Int$ dengan ``melompat''
bolak-balik antara bilangan bulat positif dan negatif. Kita juga dapat
memandang $f$ sebagai fungsi yang didefinisikan menurut kasus sebagai berikut:
\[
f(n) = \begin{cases}
  \frac{n}{2} & \text{jika $n$ genap}\\
  -\frac{n+1}{2} & \text{jika $n$ ganjil}
  \end{cases}
\]
\end{ex}

\begin{prob}
Tunjukkan bahwa jika $A$ dan $B$ adalah !!{enumerable}, maka $A \cup B$ juga
demikian.
\end{prob}

\begin{prob}
Tunjukkan dengan induksi pada $n$ bahwa jika $A_1$, $A_2$, \dots, $A_n$
semuanya !!{enumerable}, maka $A_1 \cup \dots \cup A_n$ juga demikian.
\end{prob}

\end{document}
